\scanpage{29}\mspage{27}
où $\bar{u}$ est le chemin associé à $u$
\diagram{\begin{tikzcd}[ampersand replacement=\&,column sep=1.15em]
x \arrow[r,no head,"{u}"] \& y \arrow[r,<-,"{\mathrm{id}}"] \& y \arrow[r,no head,"{\mathrm{id}}"] \& y \arrow[r,<-,"{\mathrm{id}}"] \& y
\end{tikzcd}}.
On a donc à regarder
\diagram{\begin{tikzcd}[ampersand replacement=\&,column sep=1.15em]
x \arrow[r,no head,"{u}"] \& y \arrow[r,<-,"{\mathrm{id}}"] \& y \arrow[r,no head,"{\mathrm{id}}"] \& y \arrow[r,<-,"{\mathrm{id}}"] \& y \arrow[r,no head,"{\mathrm{id}}"] \& y \arrow[r,<-,"{u}"] \& x \arrow[r,no head,"{\mathrm{id}}"] \& x \arrow[r,<-,"{\mathrm{id}}"] \& x
\end{tikzcd}}
et voir qu'il est homotope au chemin trivial
\diagram{\begin{tikzcd}[ampersand replacement=\&,column sep=1em]
x\arrow[r,"\mathrm{id}"]\&x\arrow[r,<-,"\mathrm{id}"]\&{}\&\cdots\&x\arrow[r,no head,"\mathrm{id}"]\&\unsure{X|x}\arrow[r,<-,"\mathrm{id}"]\&x
\end{tikzcd}}
(des identités partout), et item pour
\diagram{\begin{tikzcd}[ampersand replacement=\&,column sep=1.15em]
y \arrow[r,no head,"{\mathrm{id}}"] \& y \arrow[r,<-,"{u}"] \& x \arrow[r,no head,"{\mathrm{id}}"] \& x \arrow[r,<-,"{\mathrm{id}}"] \& x \arrow[r,no head,"{u}"] \& y \arrow[r,<-,"{\mathrm{id}}"] \& y \arrow[r,"{\mathrm{id}}"] \& y \arrow[r,<-,"{\mathrm{id}}"] \& y
\end{tikzcd}}
et le chemin trivial en $y$. Cette fois, on
a \inserted{encore} deux foncteurs
\diagram{\begin{tikzcd}[ampersand replacement=\&,column sep=large]
\mathord{\tikz[baseline=0pt,x=1ex,y=1ex]{\draw[line width=.09ex,line join=round] (0,0) -- (.95,1.30) -- (1.52,0) -- cycle; \draw[line width=.09ex] (.32,0) -- (1.02,1.12);}}_1 \& \unsure{\mathscr{L}|\mathcal{L}}_4\arrow[l,shift right=1,"\varphi"']\arrow[l,shift left=1,"\psi"]\\
{} \& \varepsilon\arrow[u,"\alpha"']\arrow[ul,"{\unsure{\alpha'|\alpha_1}}"']
\end{tikzcd}}
\deleted{\unsure{transfor}} compatibles avec \inserted{objets} sources et \inserted{objets} buts, et on doit
prouver qu'ils sont homotopes dans $\underline{\mathrm{Hom}}^{!}(\mathscr{L}_4,\mathord{\tikz[baseline=0pt,x=1ex,y=1ex]{\draw[line width=.09ex,line join=round] (0,0) -- (.95,1.30) -- (1.52,0) -- cycle; \draw[line width=.09ex] (.32,0) -- (1.02,1.12);}}_1)$.
Cela doit donc résulter du même lemme
général. On a \unsure{gagné|``gagné''} -- par une interprétation
simple de $X\unsure{\Sigma|Z}^{-1}$, et en particulier des
$\operatorname{Hom}_{X\Sigma^{-1}}(x,y)$, \inserted{et de leur composition,} en termes de la géométrie
des espaces de chemins :
\[
\left\{\begin{aligned}
&\operatorname{Hom}_{X\Sigma^{-1}}(x,y)\simeq\pi_0\bigl(\underline{\mathrm{Ch}}_\Sigma(X;x,y)\bigr)\\
&\text{Composition des hom dans }X\Sigma^{-1}\text{ correspond à}\\
&\qquad\text{la composition des $\Sigma$-chemins dans $X$}
\end{aligned}\right.
\]
% archive sheet number: 28
